\documentclass[a4paper,12pt]{report}

\usepackage[T2A]{fontenc}
\usepackage[utf8]{inputenc}
\usepackage[english,russian]{babel}
\usepackage{lmodern}
\DeclareFontFamily{T2A}{lmr}{}
\DeclareFontShape{T2A}{lmr}{m}{n}{<->ssub*cmr/m/n}{}
\DeclareFontShape{T2A}{lmr}{bx}{n}{<->ssub*cmr/bx/n}{}
\DeclareFontShape{T2A}{lmr}{b}{n}{<->ssub*cmr/bx/n}{}
\DeclareFontShape{T2A}{lmr}{m}{it}{<->ssub*cmr/m/it}{}
\DeclareFontShape{T2A}{lmr}{m}{sl}{<->ssub*cmr/m/sl}{}
\usepackage{graphicx}
\usepackage{setspace}
\usepackage{fancyvrb}
\usepackage{appendix}
\usepackage{listings}
\usepackage{color}
\usepackage{amssymb}
\usepackage{amsmath}
\usepackage{hyphenat}
\usepackage{subcaption}
\usepackage{hyperref}
\usepackage{float}
\usepackage{booktabs}
\usepackage{array}
\makeatletter
\renewcommand\@biblabel[1]{#1.}
% \paragraph --- отдельной строкой: положительный отступ после заголовка, текст после него --- новый абзац.
\renewcommand\paragraph{\@startsection{paragraph}{4}{\z@}{3.25ex \@plus1ex \@minus.2ex}{1ex \@plus.2ex}{\normalfont\normalsize\bfseries}}
\makeatother

\newcommand{\myrule}[1]{\rule{#1}{0.4pt}}
\newcommand{\sign}[2][~]{{\small\myrule{#2}\\[-0.7em]\makebox[#2]{\it #1}}}

% Поля
\usepackage[top=20mm, left=30mm, right=15mm, bottom=20mm, includefoot]{geometry}
\usepackage{indentfirst}
\footskip 20pt

\definecolor{gray}{rgb}{0.5,0.5,0.5}
\setcounter{secnumdepth}{4}

\bibliographystyle{gost780s}  %% стилевой файл для оформления по ГОСТу

\lstset{
	basicstyle=\footnotesize,           % the size of the fonts that are used for the code
	numbers=left,                   % where to put the line-numbers
	numberstyle=\tiny\color{gray},
	captionpos=b,                   % sets the caption-position to bottom
	breaklines=true,                % sets automatic line breaking
}

\begin{document}
		%%%%%%%%%%%%%%%%%%%%%%%%%%%%%%%
	%%%                         %%%
	%%% Начало титульного листа %%%
	
	\thispagestyle{empty}
	{ \large
		%%% Название работы %%%
		\begin{center}
			\renewcommand{\baselinestretch}{1.2}
			\Large
			{\Large Анализ зашумленных и частично наблюдаемых данных судовых навигационных систем}
		\end{center}
		\medskip%
		
		\vspace{3\bigskipamount}%
		\begin{flushright}
			\begin{minipage}{0.5\textwidth}
				\begin{flushleft}
					Гашков С.П.
				\end{flushleft}
			\end{minipage}
		\end{flushright}
		
		\vfill%
	}
	
	%%% Конец титульного листа  %%%
	%%%                         %%%
	%%%%%%%%%%%%%%%%%%%%%%%%%%%%%%%
	
		% Размер шрифта (14пт)
	\large
	
	% Межстрочный интервал (полуторный)
	\onehalfspacing
	
	\clearpage
	\renewcommand{\thesection}{\arabic{section}}
	\graphicspath{{../pict/}}
	
	\tableofcontents
	\newpage
	
	\section{Постановка задачи}
	
	
\end{document}